\documentclass[12pt, letterpaper]{article}

% --- Packages ---
\usepackage{geometry}
 \geometry{margin=1in}
\usepackage{graphicx}
\usepackage{amsmath, amssymb}
\usepackage{booktabs}
\usepackage{hyperref}
\usepackage{cite}
\usepackage{setspace}
\onehalfspacing
\usepackage{fontspec}
\usepackage{polyglossia}
\usepackage{float}
\usepackage{svg}
\usepackage{enumitem}
\setmainlanguage{vietnamese}

\usepackage{listings}
\usepackage{xcolor}

\lstset{
    language=C,
    basicstyle=\ttfamily\small,
    keywordstyle=\color{blue},
    stringstyle=\color{red},
    commentstyle=\color{green!60!black},
    breaklines=true,
    frame=single, % Tạo khung xung quanh code
    backgroundcolor=\color{gray!5}
}

\renewcommand{\lstlistingname}{Mã nguồn}

% --- Metadata ---
\title{Giải bài tập 14}
\author{Đỗ Hoàng Gia Bảo - 23020783 \\ \small Khoa Điện tử - Viễn Thông}
\date{\today}

\begin{document}

\maketitle

\begin{abstract}
\begin{center}
    Báo cáo sẽ trình bày các bài tập trong slide 33, 34 Bài-9 
\end{center}
\end{abstract}
% --- Front Matter (Roman Numerals) ---
\newpage
\tableofcontents

% --- Main Content (Arabic Numerals) ---
\newpage
\pagenumbering{arabic}
\setcounter{page}{3} % This automatically resets the counter to 1

\section{Bài tập 1}
\subsection{Đề bài}
\begin{enumerate}[label=\alph*)]
    \item Ổ cứng của máy tính có kích thước 128 GB, kích thước
            mỗi khối trên ổ đĩa là 8 KB. Nếu hệ điều hành sử dụng
            FAT thì dung lượng bộ nhớ nhỏ nhất để lưu FAT là bao
            nhiêu ? Giải thích.
    \item Hệ điều hành sử dụng i-node để quản lý các khối dữ liệu
            của tập tin. I-node của mỗi tập tin chứa số hiệu của 12
            khối trực tiếp, 1 khối gián tiếp một cấp 1, 1 khối gián tiếp
            một cấp 2, 1 khối gián tiếp một cấp 3. Kích thước mỗi
            khối trên ổ đĩa là 4 KB, số hiệu của mối khối chiếm 4
            byte.
            \begin{enumerate}[label=\roman*)]
                \item Tính kích thước lớn nhất của tập tin ?
                \item Cần bao nhiêu khối trên ổ cứng để chứa được tập 
                        tin này ?
            \end{enumerate}
            (Giả định rằng ổ cứng luôn đủ lớn để chứa được tập tin)
\end{enumerate}
\subsection{Giải bài tập}
\subsubsection{Ý a}
Vì kích thước của ổ cứng là 128 GB và kích thước của mỗi khối trên ổ đĩa là 8 KB, số khối trong ổ cứng là:
$$ \text{Số Khối} = \frac{\text{Kích thước ổ cứng}}{\text{kích thước một khối}} $$
$$ \text{Số khối} = \frac{128 \text{ GB}}{8 \text{ KB}} = 16,777,216 \text{ (khối)} $$
Số bit cần thiết để mã hóa địa chỉ tất cả số khối trên
$$ \text{Số bit} = \log_2 (16,777,216) = 24 \text{ (bit)} $$
Để có thể lưu được số bit địa chỉ trên, kích thước của mỗi entry để lưu một địa chỉ là 24 bit. 
Đồng thời số địa chỉ mà FAT cần lưu là 16,777,216. Vì vậy kích thước của FAT là:
$$ \text{Kích thước FAT} = \text{kích thước một entry} \times \text{Số entry} $$
$$ \text{Kích thước FAT} = 24 \text{ bit} \times 16,777,216 = 402,653,184 \text{ (bit)} \approx 48 \text{ (MB)} $$
Trong trường hợp theo quy chuẩn của các kiểu FAT, vì cần 24 bit để biểu diễn toàn bộ địa chỉ khối nhớ trong ổ cứng,
sử dụng chuẩn FAT32 để biểu diễn địa chỉ. Khi đó, kích thước của FAT32 với số entry trên là:
$$ \text{Kích thước FAT} = \text{kích thước một entry} \times \text{Số entry} $$
$$ \text{Kích thước FAT} = 32 \text{ bit} \times 16,777,216 = 536,870,912 \text{ (bit)} = 64 \text{ (MB)} $$
\subsubsection{Ý b}
I-node của mỗi tập tin chứa số hiệu của 12 khối trực tiếp nên kích thước của các khối trực tiếp đó là:
$$ \text{Tổng kích thước các khối trực tiếp} = \text{Số khối trực tiếp} \times \text{Kích thước một khối} $$
$$ \text{Tổng kích thước các khối trực tiếp} = 12 \times 4 \text{ (KB)} = 48 \text{ (KB)} $$
Vì khối gián tiếp lưu số hiệu mục nên số mục có thể lưu được trên mỗi khối gián tiếp là:
$$ \text{Số mục} = \frac{\text{Kích thước một khối}}{\text{Kích thước một số hiệu}} = \frac{4 \text{ (KB)}}{4 \text{ (bytes)}} = 1024 \text{ (mục)} $$
Vì các khối gián tiếp lưu số hiệu của các khối tiếp khác, số khối mỗi loại khối gián tiếp có thể lưu là:
$$ \text{Khối gián tiếp cấp 1} = \text{Số mục} = 1024 \text{ (mục)} $$
$$ \text{Khối gián tiếp cấp 2} = \text{Số mục}^2 = 1024^2 \text{ (mục)} $$
$$ \text{Khối gián tiếp cấp 3} = \text{Số mục}^3 = 1024^3 \text{ (mục)} $$
Từ các kết quả trên, tổng số khối dữ liệu là:
$$\text{Tổng số khối dữ liệu} = 12 + 1024 + 1024^2 + 1024^3 = 1,074,791,436 \text{ khối}$$
Kích thước tập tin lớn nhất khi nó sử dụng tất cả các khối. Vì vậy kích thước lớn nhất của tập tin là:
$$\text{Kích thước tập tin} = 1,074,791,436 \times 4 \text{ KB} \approx 4 \text{ (TB)}$$
\textbf{Kết luận}
\begin{enumerate}[label=\roman*)]
    \item $ \text{Kích thước lớn nhất của tập tin} \approx 4 \text{ (TB)} $
    \item $ \text{Số khối trên ổ cứng để chứa được tập tin} = 1,074,791,436 \text{ (Khối)} $
\end{enumerate}
\section{Bài tập 2}
\subsection{Đề bài}
Một tập tin được lưu trên 100 khối. Giả định rằng khối điều
khiển tập tin (và các khối chỉ mục nếu có) đều đã được
chuyển vào bộ nhớ trong. Tính số lần truy cập vào ổ đĩa để:
\begin{enumerate}[label=\alph*)]
    \item Thêm một khối vào đầu tập tin
    \item Thêm một khối vào cuối tập tin
    \item Bỏ một khối đầu tập tin
    \item Bỏ một khối cuối tập tin
\end{enumerate}
khi sử dụng phương thức cấp phát liên tục, danh sách liên
kết, và chỉ mục. Biết rằng với phương thức cấp phát liên tục,
chỉ có thể mở rộng tập tin về phía cuối tập tin. Thông tin của
các khối thêm vào được lưu ở trong bộ nhớ.
\subsection{Giải bài tập}
\begin{enumerate}
    \item Phương thức cấp phát liên tục
            \begin{enumerate}[label=\alph*)]
                \item Thêm một khối vào đầu tập tin
                
                Để thêm được khối ở đầu, cần phải tìm vùng mới gồm 101 khối và chép lại toàn bộ tập tin. Cần đọc 100 khối cũ và ghi 101 khối mới nên tổng là $100 + 101 = 201 \text{ lần}$
                \item Thêm một khối vào cuối tập tin
                
                Vì cấp phát liên tục chỉ có thể mở rộng về cuối, chỉ cần ghi thêm 1 khối mới (1 lần). Tuy nhiên, nếu khối ngay sau đó không trống thì cần tìm vùng trống liên tục mới nên sẽ tốn $100 + 101 = 201 \text{ lần}$
                \item Bỏ một khối đầu tập tin
                
                Vì dữ liệu liên tục nên phải tách 99 khối sau. $\text{Đọc 99 khối} + \text{Ghi lại 99 khối} = 198 \text{ (lần)} $
                \item Bỏ một khối cuối tập tin
                
                Vì chỉ giảm kích thước của file đi một khối nên không cần đọc ghi dữ liệu $\rightarrow 0 \text{ (lần)} $
            \end{enumerate}
    \item Phương thức cấp phát liên kết
            \begin{enumerate}[label=\alph*)]
                \item Thêm một khối vào đầu tập tin
                
                Vì phương thức sử dụng danh sách liên kết, thêm khối đầu chỉ cần trỏ đến khối đầu cũ $\rightarrow 0 \text{ (lần)} $
                \item Thêm một khối vào cuối tập tin
                
                Vì con trỏ phải cập nhật để đi về cuối. Đọc lần lượt để tìm khối cuối mất 100 lần, ghi khối cuối để cập nhật con trỏ (-1) mất 1 lần và ghi khối mới mất 1 lần $\rightarrow 100 + 1 + 1 = 102 \text{ (lần)} $
                \item Bỏ một khối đầu tập tin

                Chỉ cần đổi con trỏ đầu sang khối kế tiếp nên không cần truy cập dữ liệu trên đĩa $\rightarrow 0 \text{ (lần)} $
                \item Bỏ một khối cuối tập tin
                
                Đọc 99 khối để tìm khối áp cuối, ghi khối cuối (-1) $\rightarrow 99 + 1 = 100 \text{ (lần)} $
            \end{enumerate}
    \item Phương thức cấp phát chỉ mục
            \begin{enumerate}[label=\alph*)]
                \item Thêm một khối vào đầu tập tin
                
                Ghi khối dữ liệu mới vào bất cứ khối trống nào và cập nhật khối chỉ mục trên RAM $\rightarrow 1 \text{ (lần)}$
                \item Thêm một khối vào cuối tập tin 
                
                Ghi khối mới vào ổ đĩa. Sau đó, thêm địa chỉ của nó vào cuối danh sách trong khối chỉ mục trên RAM $\rightarrow 1 \text{ (lần)}$
                \item Bỏ một khối đầu tập tin
                
                Khối chỉ mục trên RAM xóa bỏ địa chỉ của Khối 1 $\rightarrow 0 \text{ (lần)}$
                \item Bỏ một khối cuối tập tin
                
                Khối chỉ mục trên RAM xóa bỏ địa chỉ của Khối cuối cùng $\rightarrow 0 \text{ (lần)}$
            \end{enumerate}
\end{enumerate}
\end{document}